% 第二章新增第5、6节：产业链价值分配 + 供应链关系矩阵
% 纯LaTeX片段，可被 ch02_industry_overview.tex 包含

\section{产业链结构与价值分配}

功率半导体产业链呈现典型的"微笑曲线"价值分布特征——上游设备与材料技术壁垒最高、利润率最丰厚；下游高端应用（车规级、数据中心）客户认证周期长、产品溢价显著；中游制造环节重资产投入大、规模效应显著，利润率相对偏低。理解各环节的价值分配与国产替代进度，是把握行业投资机会的基础。

\begin{exhibitbox}[图 2-4：功率半导体产业链价值分配图谱]
\centering
\vspace{0.3cm}
\begin{tikzpicture}[scale=0.92]

% 微笑曲线背景
\draw[smooth, thick, navy!40, fill=navy!8] plot coordinates {
    (0.5, 3.8) (1.5, 3.2) (2.5, 2.6) (3.5, 2.2) (4.5, 2.0)
    (5.5, 2.2) (6.5, 2.7) (7.5, 3.3) (8.5, 3.9) (9.5, 4.3)
};
\node[navy, font=\bfseries\small] at (5, 1.2) {产业链环节};
\node[navy, font=\bfseries\small, rotate=90] at (-0.3, 3) {附加值 / 毛利率};

% 上游环节
\fill[navy!25, draw=navy, thick] (0.3, 3.6) rectangle (1.7, 4.6);
\node[font=\bfseries\small, text centered] at (1, 4.25) {设备};
\node[font=\scriptsize, text centered] at (1, 3.9) {40--45\%};

\fill[navy!20, draw=navy, thick] (1.8, 3.2) rectangle (3.2, 4.2);
\node[font=\bfseries\small, text centered] at (2.5, 3.9) {衬底/材料};
\node[font=\scriptsize, text centered] at (2.5, 3.55) {25--40\%};

% 中游环节
\fill[accentblue!25, draw=accentblue, thick] (3.3, 2.3) rectangle (4.7, 3.3);
\node[font=\bfseries\small, text centered] at (4, 3.0) {制造/代工};
\node[font=\scriptsize, text centered] at (4, 2.65) {20--30\%};

\fill[accentblue!30, draw=accentblue, thick] (4.8, 2.0) rectangle (6.2, 3.0);
\node[font=\bfseries\small, text centered] at (5.5, 2.7) {封测};
\node[font=\scriptsize, text centered] at (5.5, 2.35) {15--25\%};

% 下游环节
\fill[riskgreen!25, draw=riskgreen, thick] (6.3, 2.8) rectangle (7.7, 3.8);
\node[font=\bfseries\small, text centered] at (7, 3.5) {设计 (Fabless)};
\node[font=\scriptsize, text centered] at (7, 3.15) {30--40\%};

\fill[riskgreen!30, draw=riskgreen, thick] (7.8, 3.4) rectangle (9.2, 4.4);
\node[font=\bfseries\small, text centered] at (8.5, 4.05) {车规/AI 应用};
\node[font=\scriptsize, text centered] at (8.5, 3.7) {溢价 30--100\%};

% 标注
\node[below=0.3cm of current bounding box.south, font=\scriptsize\color{steelgrey}]
    {图 2-4 功率半导体产业链微笑曲线与价值分配示意图};

\end{tikzpicture}
\vspace{0.2cm}
\sourcenote{本研究团队根据各公司年报及行业数据整理绘制}
\end{exhibitbox}

从毛利率水平看，功率半导体设备环节毛利率最高（北方华创 40--45\%），其次是功率 IC 设计（35--50\%）和 SiC 衬底（景气周期可达 40\%+）；封测环节毛利率最低（15--25\%），属于重资产、劳动密集型环节。值得注意的是，SiC 产业链的价值高度集中在上游——衬底加外延占 SiC 器件 BOM 成本的 50--60\%，这与硅基产业链价值相对分散的格局形成鲜明对比。

\begin{exhibitbox}[表 2-3：功率半导体各环节价值分配与国产替代进度]
\centering
\small
\begin{tabular}{@{}lcccccc@{}}
\toprule
\textbf{产业链环节} & \textbf{代表公司} & \textbf{毛利率 (2025A)} & \textbf{典型区间} & \textbf{价值占比} & \textbf{国产化率} & \textbf{替代阶段} \\
\midrule
\rowcolor{navy!5}
\multicolumn{6}{l}{\textit{上游：材料与设备}} \\
设备 & 北方华创 & 40--45\% & 35--45\% & 资本开支项 & 10--15\% & 差距较大 \\
SiC 衬底 & 天岳先进 & 13.1\% & 15--40\% & 30--40\% (SiC BOM) & 20--30\% (6英寸) & 技术追赶 \\
硅片 (功率级) & 立昂微 & 25--30\% & 20--35\% & 15--20\% (硅基 BOM) & 30--40\% & 快速追赶 \\
封装材料 (AMB) & 中瓷电子 & 35--40\% & 30--45\% & 5--10\% (模块 BOM) & 15--20\% & 加速突破 \\
\midrule
\rowcolor{accentblue!5}
\multicolumn{6}{l}{\textit{中游：设计与制造}} \\
功率 IC/PMIC & 晶丰明源 & 30--35\% & 35--50\% & 20--30\% & 15--20\% & 逐步突破 \\
器件设计 (Fabless) & 斯达半导 / 新洁能 & 31.6\% / 32.9\% & 30--40\% & 25--35\% & 25--40\% & 快速提升 \\
IDM (综合) & 时代电气 / 华润微 & 33.4\% / 26.4\% & 25--38\% & 40--50\% & 20--35\% & 规模替代 \\
封装测试 & 长电科技 / 通富微电 & 15--20\% & 15--25\% & 10--15\% & 40--50\% & 规模替代 \\
\midrule
\rowcolor{riskgreen!5}
\multicolumn{6}{l}{\textit{下游应用端 (产品溢价)}} \\
车规级模块 & 时代电气 / 斯达 & 溢价 30--50\% & — & — & 10--15\% & 导入验证 \\
工业级模块 & 斯达半导 / 士兰微 & 溢价 15--25\% & — & — & 30--40\% & 快速提升 \\
消费级器件 & 扬杰 / 华润微 & 基准水平 & — & — & 50--70\% & 规模替代 \\
AI/数据中心 & 士兰微 / 晶丰明源 & 溢价 20--40\% & — & — & 10--20\% & 加速突破 \\
\bottomrule
\end{tabular}
\vspace{0.2cm}
\sourcenote{各公司 2025 年年报、卖方一致预期，本研究团队整理}
\end{exhibitbox}

表 2-3 揭示了一个重要的投资规律：\textbf{国产化率越低的环节，未来成长空间越大，但短期兑现难度也越高。} 设备、SiC 衬底、车规级模块、功率 IC 等环节国产化率均低于 30\%，是国产替代的"深水区"，也是长期成长空间最大的方向。而低压 MOSFET、二极管、封测等环节国产化率已较高，增长更多依赖行业整体需求和份额的稳步提升。

\begin{keyinsight}[价值分配视角下的投资排序]
从"成长空间 + 盈利质量 + 竞争格局"三维度综合评估，我们对功率半导体各环节的投资优先级排序为：
\vspace{0.3em}
\begin{enumerate}[leftmargin=2em]
\item \textbf{第一梯队（高成长+高壁垒）}：SiC 衬底、车规 IGBT/SiC 模块、功率 IC/PMIC——国产化率低、技术壁垒高、客户粘性强，是皇冠上的明珠
\item \textbf{第二梯队（稳健成长）}：中高压 MOSFET、超结 MOS、工业级 IGBT、AMB 基板——国产替代加速，业绩兑现确定性高
\item \textbf{第三梯队（周期+主题）}：低压 MOS、二极管、封测——国产化率高，更多是周期波动和主题性机会
\end{enumerate}
\vspace{0.3em}
投资者应根据自身风险偏好和投资周期，在不同梯队间进行配置平衡。第一梯队标的弹性大但波动也大，适合长期布局；第二梯队标的业绩确定性高，适合作为核心持仓。
\end{keyinsight}

\section{供应链关系矩阵}

供应链关系是功率半导体公司竞争力的"隐性资产"——客户认证壁垒、供应商稳定性、上下游协同深度，直接决定了公司的成长天花板和盈利韧性。我们构建了覆盖上游材料设备至下游应用的全产业链供应链关系矩阵，并采用四级可信度标注体系，为投资决策提供更精细的依据。

\vspace{0.3em}

\textbf{可信度标注体系说明}：
\begin{itemize}[leftmargin=2em]
\item \textbf{\textcolor{navy}{Confirmed (C)}}：年报/公告/监管文件直接披露，或公司 IR 会议官方确认，可靠性极高
\item \textbf{\textcolor{riskgreen}{Broker-Stated (B)}}：主流卖方研报明确表述，或行业协会/咨询机构数据，可靠性较高（注意卖方偏差）
\item \textbf{\textcolor{riskamber}{Inferred (I)}}：基于行业逻辑、产业链关系的合理推断，需持续验证，可靠性中等
\item \textbf{\textcolor{riskred}{Market-Rumor (M)}}：市场传闻、自媒体报道，不可作为决策依据，可靠性低
\end{itemize}

\begin{exhibitbox}[表 2-4：功率半导体核心公司供应链关系矩阵]
\centering
\scriptsize
\renewcommand{\arraystretch}{1.2}
\begin{tabular}{@{}lllccl@{}}
\toprule
\textbf{环节} & \textbf{公司} & \textbf{供应链关系} & \textbf{方向} & \textbf{可信度} & \textbf{核心依据} \\
\midrule
\rowcolor{navy!5}
\multicolumn{6}{l}{\textit{上游材料设备 $\rightarrow$ 中游制造}} \\
SiC 衬底 & 天岳先进 & 博世 (Bosch) & 下游 & \textcolor{navy}{\textbf{C}} & 年报 + 卓越供应商奖 \\
& & 意法半导体 (ST) & 下游 & \textcolor{riskgreen}{\textbf{B}} & 华泰/华源研报 \\
& & 富士康 (鸿海) & 下游 & \textcolor{navy}{\textbf{C}} & 2026Q1 业绩说明会 \\
& & 时代电气 / 士兰微 / 三安 & 下游 & \textcolor{riskgreen}{\textbf{B}} & 招商/华泰研报 \\
硅片 & 立昂微 & 华润微 / 士兰微 / 斯达 & 下游 & \textcolor{riskgreen}{\textbf{B}} & 海通研报 \\
& 沪硅产业 & 中芯国际 / 华虹 / 华润微 & 下游 & \textcolor{riskgreen}{\textbf{B}} & 海通研报 \\
设备 & 北方华创 & 华润微 / 士兰微 / 时代电气 & 下游 & \textcolor{riskgreen}{\textbf{B}} & 海通研报 \\
& 中微公司 & 三安光电 / 华润微 & 下游 & \textcolor{riskgreen}{\textbf{B}} & 海通研报 \\
封装材料 & 中瓷电子 & 时代电气 / 斯达 / 比亚迪半导 & 下游 & \textcolor{riskgreen}{\textbf{B}} & 广发研报 \\
\midrule
\rowcolor{accentblue!5}
\multicolumn{6}{l}{\textit{中游制造 $\rightarrow$ 下游应用 (IDM 模式)}} \\
\multirow{3}{*}{IGBT+SiC IDM} & \multirow{3}{*}{时代电气} & 大众 / 丰田 / 通用 / 小米 / 奇瑞 & 汽车 & \textcolor{navy}{\textbf{C}} & 年报 + IR 交流会 \\
& & 中国中车 / 国铁集团 & 轨交 & \textcolor{navy}{\textbf{C}} & 年报 \\
& & 阳光电源 / 华为数字能源 & 光储 & \textcolor{riskgreen}{\textbf{B}} & 中信建投研报 \\
\multirow{3}{*}{全品类 IDM} & \multirow{3}{*}{华润微} & 国内一线消费品牌 & 消费 & \textcolor{navy}{\textbf{C}} & 年报 \\
& & 光模块头部客户 / 服务器电源厂 & AI & \textcolor{navy}{\textbf{C}} & 2026.4 业绩说明会 \\
& & 车规级验证中 & 汽车 & \textcolor{riskgreen}{\textbf{B}} & 国信研报 \\
\multirow{2}{*}{IDM 扩产} & \multirow{2}{*}{士兰微} & 美的 / 格力 / 海尔 & 消费 & \textcolor{riskgreen}{\textbf{B}} & 国信研报 \\
& & AI 大客户 (Dr.MOS/eFuse) & AI & \textcolor{navy}{\textbf{C}} & 2025.11 业绩说明会 \\
\multirow{3}{*}{分立 IDM} & \multirow{3}{*}{扬杰科技} & 博世 / 大陆集团 (Tier1) & 汽车 & \textcolor{navy}{\textbf{C}} & 年报 + 2026.5 电话会 \\
& & AI 头部企业 (联合开发) & AI & \textcolor{navy}{\textbf{C}} & 2026.5 电话会 \\
& & 光伏 / 储能 / 充电桩 & 能源 & \textcolor{riskgreen}{\textbf{B}} & 国金研报 \\
化合物 IDM & 三安光电 & 手机 PA 厂商 / 全球照明厂 & RF/LED & \textcolor{navy}{\textbf{C}} & 年报 \\
\midrule
\rowcolor{accentblue!5}
\multicolumn{6}{l}{\textit{中游设计 $\rightarrow$ 下游应用 (Fabless 模式)}} \\
IGBT 设计 & 斯达半导 & 多家头部车企 (主驱/OBC) & 汽车 & \textcolor{navy}{\textbf{C}} & 年报 \\
& & 汇川技术 / 英威腾 & 工控 & \textcolor{riskgreen}{\textbf{B}} & 申万/国泰君安研报 \\
MOS 设计 & 新洁能 & 消费电子 / 工业电源 & 多领域 & \textcolor{riskgreen}{\textbf{B}} & 中信/中金研报 \\
超结 MOS & 东微半导 & 充电桩 / 光伏 / 服务器 PSU & 能源+AI & \textcolor{riskgreen}{\textbf{B}} & 中信/中金研报 \\
电源管理 IC & 晶丰明源 & 照明龙头 / 家电品牌 & 消费 & \textcolor{navy}{\textbf{C}} & 年报 \\
\midrule
\rowcolor{riskgreen!5}
\multicolumn{6}{l}{\textit{下游应用领域国产化率一览}} \\
\multicolumn{2}{l}{新能源车主驱 (IGBT/SiC)} & 时代电气 / 斯达 / 比亚迪半导 & — & \multicolumn{2}{c}{国产化率 10--15\%} \\
\multicolumn{2}{l}{光伏逆变器 (IGBT/MOS)} & 斯达 / 时代 / 扬杰 & — & \multicolumn{2}{c}{国产化率 30--40\%} \\
\multicolumn{2}{l}{AI 服务器 (Dr.MOS/PMIC)} & 士兰微 / 华润微 / 晶丰明源 & — & \multicolumn{2}{c}{国产化率 10--20\%} \\
\multicolumn{2}{l}{工业控制 (变频器/伺服)} & 斯达 / 时代 / 士兰微 & — & \multicolumn{2}{c}{国产化率 30--40\%} \\
\multicolumn{2}{l}{消费电子快充 (GaN/MOS)} & 三安 / 晶丰明源 / 英诺赛科 & — & \multicolumn{2}{c}{国产化率 35--45\%} \\
\bottomrule
\end{tabular}
\vspace{0.2cm}
\sourcenote{各公司年报、业绩说明会、卖方研报（中信/华泰/海通/中金等），本研究团队整理}
\end{exhibitbox}

表 2-4 呈现的供应链关系矩阵，核心价值在于两个维度的洞察。第一，\textbf{客户质量决定成长天花板}——进入博世、大众、华为等头部客户供应链的公司，其成长确定性远高于服务中小客户的公司。例如时代电气凭借大众、丰田、通用三大海外车企定点，车规业务的长期成长路径最为清晰；扬杰科技进入博世、大陆等国际 Tier1 供应链，海外收入占比持续提升。第二，\textbf{供应链多元化决定抗风险能力}——单一客户或单一领域依赖度高的公司，面临更大的业绩波动风险。天岳先进、东微半导等公司客户集中度相对较高，需持续关注客户结构的多元化进展。

\begin{exhibitbox}[图 2-5：核心功率半导体公司供应链定位图谱]
\centering
\vspace{0.3cm}
\begin{tikzpicture}[scale=0.95]

% 坐标轴
\draw[->, thick, steelgrey] (0,0) -- (12,0) node[right, font=\small] {下游客户覆盖广度 $\rightarrow$};
\draw[->, thick, steelgrey] (0,0) -- (0,7) node[above, font=\small] {$\uparrow$ 客户质量/认证等级};

% 网格线
\draw[dashed, gray!20] (0,2) -- (12,2);
\draw[dashed, gray!20] (0,4) -- (12,4);
\draw[dashed, gray!20] (0,6) -- (12,6);
\draw[dashed, gray!20] (3,0) -- (3,7);
\draw[dashed, gray!20] (6,0) -- (6,7);
\draw[dashed, gray!20] (9,0) -- (9,7);

% Y轴标签
\node[left, font=\scriptsize, steelgrey] at (0,1) {消费级};
\node[left, font=\scriptsize, steelgrey] at (0,3) {工业级};
\node[left, font=\scriptsize, steelgrey] at (0,5) {车规级};
\node[left, font=\scriptsize, steelgrey] at (0,6.5) {国际 Tier1};

% X轴标签
\node[below, font=\scriptsize, steelgrey] at (1.5,0) {单一领域};
\node[below, font=\scriptsize, steelgrey] at (4.5,0) {2--3个领域};
\node[below, font=\scriptsize, steelgrey] at (7.5,0) {多领域覆盖};
\node[below, font=\scriptsize, steelgrey] at (10.5,0) {全领域布局};

% 象限标注
\node[navy, font=\bfseries\footnotesize, above right] at (6.2, 5.8) {黄金象限：车规级 + 广覆盖};
\node[steelgrey, font=\footnotesize, above right] at (0.2, 5.8) {专精型};
\node[steelgrey, font=\footnotesize, above right] at (0.2, 0.8) {细分赛道型};
\node[steelgrey, font=\footnotesize, above right] at (9.2, 0.8) {平台型 (消费为主)};

% 公司点位 - 第一梯队（黄金象限附近）
\node[draw=navy, fill=navy!20, circle, minimum size=12pt, inner sep=0pt] (sddz) at (5, 6.1) {};
\node[below=2pt of sddz, font=\scriptsize, navy] {\textbf{时代电气}};

\node[draw=navy, fill=navy!15, circle, minimum size=10pt, inner sep=0pt] (std) at (4, 5.2) {};
\node[below=2pt of std, font=\scriptsize] {斯达半导};

\node[draw=navy, fill=navy!15, circle, minimum size=10pt, inner sep=0pt] (yj) at (7, 5.0) {};
\node[below=2pt of yj, font=\scriptsize] {扬杰科技};

% 第二梯队
\node[draw=accentblue, fill=accentblue!20, circle, minimum size=12pt, inner sep=0pt] (hrw) at (8, 4.2) {};
\node[below=2pt of hrw, font=\scriptsize, accentblue] {\textbf{华润微}};

\node[draw=accentblue, fill=accentblue!20, circle, minimum size=12pt, inner sep=0pt] (slw) at (7, 3.8) {};
\node[below=2pt of slw, font=\scriptsize, accentblue] {\textbf{士兰微}};

\node[draw=accentblue, fill=accentblue!15, circle, minimum size=10pt, inner sep=0pt] (xjn) at (5, 3.0) {};
\node[below=2pt of xjn, font=\scriptsize] {新洁能};

\node[draw=accentblue, fill=accentblue!15, circle, minimum size=10pt, inner sep=0pt] (dw) at (4, 2.6) {};
\node[below=2pt of dw, font=\scriptsize] {东微半导};

% 材料设备公司
\node[draw=riskamber, fill=riskamber!20, circle, minimum size=10pt, inner sep=0pt] (ty) at (2.5, 5.8) {};
\node[below=2pt of ty, font=\scriptsize, riskamber] {天岳先进};

\node[draw=riskamber, fill=riskamber!15, circle, minimum size=10pt, inner sep=0pt] (san) at (9, 3.0) {};
\node[below=2pt of san, font=\scriptsize] {三安光电};

% 设计公司
\node[draw=riskgreen, fill=riskgreen!15, circle, minimum size=9pt, inner sep=0pt] (jf) at (6, 1.8) {};
\node[below=2pt of jf, font=\scriptsize] {晶丰明源};

% 图例
\node[below right, font=\scriptsize] at (0.5, 0.5) {\textbf{图例：}};
\node[draw=navy, fill=navy!20, circle, minimum size=8pt, inner sep=0pt] at (2, 0.5) {};
\node[right, font=\scriptsize] at (2.5, 0.5) {IDM 龙头};
\node[draw=accentblue, fill=accentblue!20, circle, minimum size=8pt, inner sep=0pt] at (4, 0.5) {};
\node[right, font=\scriptsize] at (4.5, 0.5) {多品类厂商};
\node[draw=riskamber, fill=riskamber!20, circle, minimum size=8pt, inner sep=0pt] at (6.5, 0.5) {};
\node[right, font=\scriptsize] at (7, 0.5) {材料/化合物};
\node[draw=riskgreen, fill=riskgreen!20, circle, minimum size=8pt, inner sep=0pt] at (9.5, 0.5) {};
\node[right, font=\scriptsize] at (10, 0.5) {设计公司};

\node[below=0.5cm of current bounding box.south, font=\scriptsize\color{steelgrey}]
    {图 2-5 核心公司供应链定位示意（基于客户质量与覆盖广度二维评估）};

\end{tikzpicture}
\vspace{0.2cm}
\sourcenote{本研究团队基于供应链关系矩阵综合评估绘制}
\end{exhibitbox}

图 2-5 将核心公司置于"客户质量"与"覆盖广度"的二维坐标系中，直观呈现了各公司的供应链定位。\textbf{时代电气}凭借车规级客户（大众、丰田、通用等）和轨交领域的深厚积累，位于图谱的"黄金象限"——客户质量最高且覆盖多领域。\textbf{华润微}和\textbf{士兰微}作为 IDM 平台型公司，客户覆盖广度领先，但车规级客户的突破仍在进行中。\textbf{天岳先进}客户质量很高（博世、意法等国际一线），但下游应用领域相对集中（以 SiC 衬底间接下游为主）。

\begin{keyinsight}[供应链视角下的选股逻辑]
从供应链关系质量出发，我们提炼出三条选股准则：
\vspace{0.3em}
\begin{enumerate}[leftmargin=2em]
\item \textbf{优先选择已进入国际 Tier1/头部车企供应链的公司}——客户认证是最高门槛，进入意味着长期成长确定性。时代电气（大众/丰田/通用）、扬杰科技（博世/大陆）、天岳先进（博世/意法）是典型代表
\item \textbf{关注客户结构多元化的平台型公司}——单一领域依赖度低，抗周期能力更强。华润微、士兰微在消费、工业、汽车、AI 多领域布局，业绩平滑能力优于单一赛道公司
\item \textbf{警惕"故事型"供应链标的}——仅停留在"研发中""验证中"阶段、缺乏 confirmed 级别客户关系的公司，业绩兑现风险较高。应以 L1/L2 级数据为决策依据，L3 级以下仅作参考
\end{enumerate}
\vspace{0.3em}
我们建议将供应链关系质量作为功率半导体选股的"一票否决"指标之一——缺乏高质量客户验证的公司，即使技术故事再动听，也应谨慎参与。
\end{keyinsight}

\vspace{0.3em}

\textbf{数据局限性提示}：本供应链关系矩阵基于公开信息整理，存在三方面局限。其一，A 股上市公司不强制披露具体客户名称和收入金额，客户占比数据多基于券商研报估算，精确性有限；其二，供应链关系处于动态变化中，客户认证、定点进展随时间持续更新，本矩阵仅反映 2026 年 6 月时点的大致格局；其三，部分公司（如斯达半导、新洁能）公开 IR 信息透明度较低，相关分析更多依赖卖方研报，信息颗粒度较粗。投资者在使用本矩阵时，应结合最新公告和调研信息持续验证。
